% --- Archon physics formalization source begin ---
% archon:physics
% archon:covers PhyXMiniProblems/problem_phyx_mini_0661.lean
% archon:source-report reports/phyx_mini/problem_phyx_mini_0661.source.json

\chapter{Physics problem phyx\_mini\_0661}
\label{ch:PhyXMiniProblems_problem_phyx_mini_0661}

\paragraph{Problem source.}
\#\# Physical scenario

A valve in the cylinder shown in figure has a cross-sectional area of \$11 \textbackslash{}, \textbackslash{}text\{cm\}\textasciicircum{}2\$ with a pressure of \$735 \textbackslash{}, \textbackslash{}text\{kPa\}\$ inside the cylinder and \$99 \textbackslash{}, \textbackslash{}text\{kPa\}\$ outside.

\#\# Question

How large a force is needed to open the valve?

\#\# Answer choices

- A: A: 836 N

- B: B: 70 N

- C: C: 634 N

- D: D: 700 N

\#\# Figure caption (auxiliary; use the image as primary evidence)

The image depicts a diagram of a piston-cylinder system with a valve. 

Key components:

1. **Cylinder**: This contains fluid, shown as a light blue area labeled with the pressure \textbackslash{}( P\_\{\textbackslash{}text\{cyl\}\} \textbackslash{}).

2. **Piston**: A solid part within the cylinder, positioned at the bottom of the light blue area. It can move up or down to change the volume inside the cylinder.

3. **Valve**: Situated at the top of the cylinder, connecting to a curved pipe leading outside. The valve has an area labeled \textbackslash{}( A\_\{\textbackslash{}text\{valve\}\} \textbackslash{}).

4. **Pipes**: Connected to the valve, indicating an entry or exit point for the fluid. The outside pressure is labeled \textbackslash{}( P\_\{\textbackslash{}text\{outside\}\} \textbackslash{}).

The relationships shown display how fluid flow and pressures are interrelated in such a thermodynamic system.

\#\# Dataset metadata

Statics; Physical Model Grounding Reasoning, Multi-Formula Reasoning

\paragraph{Recorded answer/context.}
D: D: 700 N

\paragraph{Figure/image path.}
/root/proposal\_for\_physic/science-mango/phyx\_mini\_run/phyx\_data/test\_image/661.png

\paragraph{Formalization target.}
create a compiling Lean file with sorry bodies at `PhyXMiniProblems/problem\_phyx\_mini\_0661.lean`. The Lean declarations must preserve the physical quantities, dimensions or dimensional roles, figure labels, governing-law hypotheses, and final relation expressed by this problem.
Use Mathlib/Physlib names found through LeanExplore where available. If a physics API is missing, introduce faithful local abstractions rather than scalar placeholder aliases.

\begin{theorem}[Physics formalization target]
\label{thm:physics:phyx_mini_0661:target}
\lean{PhyXMiniProblems.ProblemPhyXMini0661.problem_phyx_mini_0661}
\uses{def:physics:phyx-mini-0661:phyxminiproblems-problemphyxmini0661-forceinnewtons, def:physics:phyx-mini-0661:phyxminiproblems-problemphyxmini0661-valveopeningsetup, def:physics:phyx-mini-0661:phyxminiproblems-problemphyxmini0661-matchesvalvecylinderscenario, def:physics:phyx-mini-0661:phyxminiproblems-problemphyxmini0661-matchessuppliedvalvefigure, def:physics:phyx-mini-0661:phyxminiproblems-problemphyxmini0661-matchesproblemreadouts, def:physics:phyx-mini-0661:phyxminiproblems-problemphyxmini0661-satisfiesvalveopeningstatics, def:physics:phyx-mini-0661:phyxminiproblems-problemphyxmini0661-answerchoice, def:physics:phyx-mini-0661:phyxminiproblems-problemphyxmini0661-displayedforceinnewtons, def:physics:phyx-mini-0661:phyxminiproblems-problemphyxmini0661-roundstonearestnewton, def:physics:phyx-mini-0661:phyxminiproblems-problemphyxmini0661-isnearestdisplayedforce}
The assigned autoformalize agent should translate this physics problem into Lean declarations in the covered file, with theorem and lemma proof bodies written as `by sorry`.
\end{theorem}
\begin{proof}
This is an autoformalization task, not a proof task. Produce statements that can later be proved by the physics prover without weakening the physical model.
\end{proof}

% --- Archon declaration topology begin ---
\section{Declaration topology}
\begin{definition}[Force Magnitude]
  \label{def:physics:phyx-mini-0661:phyxminiproblems-problemphyxmini0661-forcemagnitude}
  \lean{PhyXMiniProblems.ProblemPhyXMini0661.ForceMagnitude}
  A nonnegative physical force magnitude, independent of the chosen units
\end{definition}
\begin{proof}
  This is the declared model definition of Force Magnitude.
\end{proof}
\begin{definition}[area In Square Meters]
  \label{def:physics:phyx-mini-0661:phyxminiproblems-problemphyxmini0661-areainsquaremeters}
  \lean{PhyXMiniProblems.ProblemPhyXMini0661.areaInSquareMeters}
  Coherent-SI square-metre readout of a physical area
\end{definition}
\begin{proof}
  This is the declared model definition of area In Square Meters.
\end{proof}
\begin{definition}[area In Square Centimeters]
  \label{def:physics:phyx-mini-0661:phyxminiproblems-problemphyxmini0661-areainsquarecentimeters}
  \lean{PhyXMiniProblems.ProblemPhyXMini0661.areaInSquareCentimeters}
  \uses{def:physics:phyx-mini-0661:phyxminiproblems-problemphyxmini0661-areainsquaremeters}
  Square-centimetre readout used for the valve-area datum
\end{definition}
\begin{proof}
  This is the declared model definition of area In Square Centimeters.
\end{proof}
\begin{definition}[pressure In Pascals]
  \label{def:physics:phyx-mini-0661:phyxminiproblems-problemphyxmini0661-pressureinpascals}
  \lean{PhyXMiniProblems.ProblemPhyXMini0661.pressureInPascals}
  Coherent-SI pascal readout of a physical pressure
\end{definition}
\begin{proof}
  This is the declared model definition of pressure In Pascals.
\end{proof}
\begin{definition}[pressure In Kilopascals]
  \label{def:physics:phyx-mini-0661:phyxminiproblems-problemphyxmini0661-pressureinkilopascals}
  \lean{PhyXMiniProblems.ProblemPhyXMini0661.pressureInKilopascals}
  \uses{def:physics:phyx-mini-0661:phyxminiproblems-problemphyxmini0661-pressureinpascals}
  Kilopascal readout used for both pressure data in the problem
\end{definition}
\begin{proof}
  This is the declared model definition of pressure In Kilopascals.
\end{proof}
\begin{definition}[force In Newtons]
  \label{def:physics:phyx-mini-0661:phyxminiproblems-problemphyxmini0661-forceinnewtons}
  \lean{PhyXMiniProblems.ProblemPhyXMini0661.forceInNewtons}
  \uses{def:physics:phyx-mini-0661:phyxminiproblems-problemphyxmini0661-forcemagnitude}
  Coherent-SI newton readout of a physical force magnitude
\end{definition}
\begin{proof}
  This is the declared model definition of force In Newtons.
\end{proof}
\begin{definition}[pressure Force In Newtons]
  \label{def:physics:phyx-mini-0661:phyxminiproblems-problemphyxmini0661-pressureforceinnewtons}
  \lean{PhyXMiniProblems.ProblemPhyXMini0661.pressureForceInNewtons}
  \uses{def:physics:phyx-mini-0661:phyxminiproblems-problemphyxmini0661-areainsquaremeters, def:physics:phyx-mini-0661:phyxminiproblems-problemphyxmini0661-pressureinpascals}
  The normal-force magnitude exerted by a uniform pressure on a planar area, using the coherent SI identity 1 Pa * 1 m² = 1 N
\end{definition}
\begin{proof}
  This is the declared model definition of pressure Force In Newtons.
\end{proof}
\begin{definition}[Figure Label]
  \label{def:physics:phyx-mini-0661:phyxminiproblems-problemphyxmini0661-figurelabel}
  \lean{PhyXMiniProblems.ProblemPhyXMini0661.FigureLabel}
  Literal symbolic labels visible in the supplied raster 661.png
\end{definition}
\begin{proof}
  This is the declared model definition of Figure Label.
\end{proof}
\begin{definition}[Figure Region]
  \label{def:physics:phyx-mini-0661:phyxminiproblems-problemphyxmini0661-figureregion}
  \lean{PhyXMiniProblems.ProblemPhyXMini0661.FigureRegion}
  Regions or components to which the visible labels are attached
\end{definition}
\begin{proof}
  This is the declared model definition of Figure Region.
\end{proof}
\begin{definition}[Valve Figure]
  \label{def:physics:phyx-mini-0661:phyxminiproblems-problemphyxmini0661-valvefigure}
  \lean{PhyXMiniProblems.ProblemPhyXMini0661.ValveFigure}
  \uses{def:physics:phyx-mini-0661:phyxminiproblems-problemphyxmini0661-figurelabel, def:physics:phyx-mini-0661:phyxminiproblems-problemphyxmini0661-figureregion}
  Qualitative information represented by the piston-cylinder diagram
\end{definition}
\begin{proof}
  This is the declared model definition of Valve Figure.
\end{proof}
\begin{definition}[Valve Opening Setup]
  \label{def:physics:phyx-mini-0661:phyxminiproblems-problemphyxmini0661-valveopeningsetup}
  \lean{PhyXMiniProblems.ProblemPhyXMini0661.ValveOpeningSetup}
  \uses{def:physics:phyx-mini-0661:phyxminiproblems-problemphyxmini0661-forcemagnitude, def:physics:phyx-mini-0661:phyxminiproblems-problemphyxmini0661-valvefigure}
  The independent physical quantities of the problem. In particular, requiredOpeningForce is not defined from a displayed answer choice
\end{definition}
\begin{proof}
  This is the declared model definition of Valve Opening Setup.
\end{proof}
\begin{definition}[Matches Valve Cylinder Scenario]
  \label{def:physics:phyx-mini-0661:phyxminiproblems-problemphyxmini0661-matchesvalvecylinderscenario}
  \lean{PhyXMiniProblems.ProblemPhyXMini0661.MatchesValveCylinderScenario}
  \uses{def:physics:phyx-mini-0661:phyxminiproblems-problemphyxmini0661-valveopeningsetup}
  The two stated pressures are uniform loads on opposite sides of the valve
\end{definition}
\begin{proof}
  This is the declared model definition of Matches Valve Cylinder Scenario.
\end{proof}
\begin{definition}[Matches Supplied Valve Figure]
  \label{def:physics:phyx-mini-0661:phyxminiproblems-problemphyxmini0661-matchessuppliedvalvefigure}
  \lean{PhyXMiniProblems.ProblemPhyXMini0661.MatchesSuppliedValveFigure}
  \uses{def:physics:phyx-mini-0661:phyxminiproblems-problemphyxmini0661-valveopeningsetup}
  Primary-image evidence. This records the three labels and the qualitative piston, valve, stem, and pipe geometry without imposing a force value
\end{definition}
\begin{proof}
  This is the declared model definition of Matches Supplied Valve Figure.
\end{proof}
\begin{definition}[Matches Problem Readouts]
  \label{def:physics:phyx-mini-0661:phyxminiproblems-problemphyxmini0661-matchesproblemreadouts}
  \lean{PhyXMiniProblems.ProblemPhyXMini0661.MatchesProblemReadouts}
  \uses{def:physics:phyx-mini-0661:phyxminiproblems-problemphyxmini0661-areainsquarecentimeters, def:physics:phyx-mini-0661:phyxminiproblems-problemphyxmini0661-pressureinkilopascals, def:physics:phyx-mini-0661:phyxminiproblems-problemphyxmini0661-valveopeningsetup}
  Numerical values stated in the problem. This contains no opening-force readout and selects no answer choice
\end{definition}
\begin{proof}
  This is the declared model definition of Matches Problem Readouts.
\end{proof}
\begin{definition}[Satisfies Valve Opening Statics]
  \label{def:physics:phyx-mini-0661:phyxminiproblems-problemphyxmini0661-satisfiesvalveopeningstatics}
  \lean{PhyXMiniProblems.ProblemPhyXMini0661.SatisfiesValveOpeningStatics}
  \uses{def:physics:phyx-mini-0661:phyxminiproblems-problemphyxmini0661-forceinnewtons, def:physics:phyx-mini-0661:phyxminiproblems-problemphyxmini0661-pressureforceinnewtons, def:physics:phyx-mini-0661:phyxminiproblems-problemphyxmini0661-valveopeningsetup}
  Quasistatic normal-force balance at the threshold of opening. The pressure load from the cylinder side is balanced by the opposing outside-pressure load and the applied opening-force magnitude. This is the governing statics law; it contains no problem-specific numerical force and no answer label
\end{definition}
\begin{proof}
  This is the declared model definition of Satisfies Valve Opening Statics.
\end{proof}
\begin{definition}[Answer Choice]
  \label{def:physics:phyx-mini-0661:phyxminiproblems-problemphyxmini0661-answerchoice}
  \lean{PhyXMiniProblems.ProblemPhyXMini0661.AnswerChoice}
  Labels of the four answers printed with the problem
\end{definition}
\begin{proof}
  This is the declared model definition of Answer Choice.
\end{proof}
\begin{definition}[displayed Force In Newtons]
  \label{def:physics:phyx-mini-0661:phyxminiproblems-problemphyxmini0661-displayedforceinnewtons}
  \lean{PhyXMiniProblems.ProblemPhyXMini0661.displayedForceInNewtons}
  \uses{def:physics:phyx-mini-0661:phyxminiproblems-problemphyxmini0661-answerchoice}
  Newton values printed beside the four answer labels. This records the candidate list without asserting which candidate follows from the model
\end{definition}
\begin{proof}
  This is the declared model definition of displayed Force In Newtons.
\end{proof}
\begin{definition}[Rounds To Nearest Newton]
  \label{def:physics:phyx-mini-0661:phyxminiproblems-problemphyxmini0661-roundstonearestnewton}
  \lean{PhyXMiniProblems.ProblemPhyXMini0661.RoundsToNearestNewton}
  \uses{def:physics:phyx-mini-0661:phyxminiproblems-problemphyxmini0661-forcemagnitude, def:physics:phyx-mini-0661:phyxminiproblems-problemphyxmini0661-forceinnewtons}
  A force readout rounds to a displayed whole-newton value
\end{definition}
\begin{proof}
  This is the declared model definition of Rounds To Nearest Newton.
\end{proof}
\begin{definition}[Is Nearest Displayed Force]
  \label{def:physics:phyx-mini-0661:phyxminiproblems-problemphyxmini0661-isnearestdisplayedforce}
  \lean{PhyXMiniProblems.ProblemPhyXMini0661.IsNearestDisplayedForce}
  \uses{def:physics:phyx-mini-0661:phyxminiproblems-problemphyxmini0661-forceinnewtons, def:physics:phyx-mini-0661:phyxminiproblems-problemphyxmini0661-valveopeningsetup, def:physics:phyx-mini-0661:phyxminiproblems-problemphyxmini0661-answerchoice, def:physics:phyx-mini-0661:phyxminiproblems-problemphyxmini0661-displayedforceinnewtons}
  A displayed force choice is at least as close as every alternative
\end{definition}
\begin{proof}
  This is the declared model definition of Is Nearest Displayed Force.
\end{proof}
% --- Archon declaration topology end ---
% --- Archon physics formalization source end ---
